\documentclass[10pt]{extarticle}
\usepackage[a4paper,margin=1.5cm]{geometry}
\usepackage[T1]{fontenc}
\usepackage{amsmath,amssymb}
\usepackage{graphicx}
\usepackage[export]{adjustbox}
\usepackage{multicol}
\usepackage{enumitem}
\usepackage{booktabs,array,tabularx}
\setlength{\columnsep}{0.7cm}
\setlength{\parindent}{0pt}
\setlength{\parskip}{3pt}
\pagestyle{empty}
\makeatletter
% figure registry (figdefs.tex is generated by the extract step)
\newcommand{\qfigdef}[4]{\@namedef{qf@file@#1}{#2}\@namedef{qf@w@#1}{#3}\@namedef{qf@wide@#1}{#4}}
\newcommand{\qf@check}[1]{\@ifundefined{qf@file@#1}{\PackageError{score2}{Unknown figure id #1}{}}{}}
\newcommand{\qf@img}[4][]{\includegraphics[width=#2\dimexpr\@nameuse{qf@w@#4}pt\relax,max width=#3,max height=0.45\textheight,#1]{\@nameuse{qf@file@#4}}}
% Full-width questions: finalize wraps every question that contains a wide figure in \qfullbegin...\qfullend,
% so the WHOLE question (stem, figures, options) is set across the page in reading order, then the two columns
% resume. The question's head - its number, stem and first figure (\qfullsplit marks the end of the head) -
% is kept on one page, so a number is never left at a page foot with its figure overleaf; the rest may break
% across pages normally, so no big gaps are left. Without a \qfullsplit, the whole question is kept together
% if it is under 45% of a page. Never write these yourself.
\newif\ifqf@full
\newif\ifqf@boxed
\newsavebox{\qf@box}
\newcommand{\qf@open}{\global\qf@fulltrue\par\global\qf@boxedtrue
  \setbox\qf@box\vbox\bgroup\hsize\textwidth\linewidth\textwidth}
\newcommand{\qf@emit}[1]{\par\egroup\global\qf@boxedfalse
  \ifdim\dimexpr\ht\qf@box+\dp\qf@box\relax<#1\textheight\noindent\usebox{\qf@box}\par\else\unvbox\qf@box\fi}
\newcommand{\qfullsplit}{\ifqf@boxed\qf@emit{0.85}\fi}
\newcommand{\qf@close}{\ifqf@boxed\qf@emit{0.45}\fi\global\qf@fullfalse}
\newcommand{\qfullbegin}{\end{multicols}\qf@open}
\newcommand{\qfullend}{\qf@close\begin{multicols}{2}}
\newcommand{\qfullstart}{\qf@open}
\newcommand{\qfulldone}{\qf@close}
% \qfig[scale]{id}: block figure, centred, directly where it is written (size set by the tool).
% Fallback if a wide figure is not inside a full-width question: break out of the columns in place - no floats.
\newcommand{\qfig}[2][1]{\qf@check{#2}%
  \ifqf@full\qf@blk{#1}{#2}\else
  \ifnum\@nameuse{qf@wide@#2}=1 \expandafter\@firstoftwo\else\expandafter\@secondoftwo\fi
  {\end{multicols}\begin{center}\qf@img{#1}{\textwidth}{#2}\end{center}\begin{multicols}{2}}%
  {\qf@blk{#1}{#2}}\fi}
\newcommand{\qf@blk}[2]{\par\smallskip{\centering\qf@img{#1}{\linewidth}{#2}\par}\smallskip}
% \qfigi[scale]{id}: inline figure (inside option lists / tables / minipages); never breaks columns.
\newcommand{\qfigi}[2][1]{\qf@check{#2}\qf@img[valign=c]{#1}{\linewidth}{#2}}
\makeatother
% structure macros - lint reads these to check completeness and figure placement
\newcommand{\papersec}[1]{\par\bigskip{\centering\Large\bfseries PAPER-#1\par}\medskip}
\newcommand{\partsec}[2]{\par\medskip\noindent{\large\bfseries PART-#1: #2}\par\smallskip}
\newcommand{\sect}[1]{\par\smallskip\noindent\textbf{#1}\par}
\newcommand{\qn}[1]{\par\smallskip\noindent\textbf{#1.}~}
\newcommand{\paraq}[2]{\par\smallskip\noindent\textbf{Paragraph for Questions #1 and #2}\par}
\newenvironment{qoptions}{\begin{itemize}[leftmargin=1.8em,itemsep=1pt,topsep=2pt,label={}]}{\end{itemize}}
\newcommand{\opt}[2]{\item[(#1)] #2}
% figure registry (generated by score2_tools.py finalize)
\qfigdef{p1_phy_s0_q03_1}{SRGBT - Test 06 (29 Aug 2021)-img/p1_phy_s0_q03_1.png}{40.1}{0}
\qfigdef{p1_phy_s0_q04_1}{SRGBT - Test 06 (29 Aug 2021)-img/p1_phy_s0_q04_1.png}{74.8}{0}
\qfigdef{p1_phy_s0_q05_1}{SRGBT - Test 06 (29 Aug 2021)-img/p1_phy_s0_q05_1.png}{72.3}{0}
\qfigdef{p1_phy_s0_q06_1}{SRGBT - Test 06 (29 Aug 2021)-img/p1_phy_s0_q06_1.png}{73.2}{0}
\qfigdef{p1_phy_s0_q06_2}{SRGBT - Test 06 (29 Aug 2021)-img/p1_phy_s0_q06_2.png}{82.8}{0}
\qfigdef{p1_phy_s0_q06_3}{SRGBT - Test 06 (29 Aug 2021)-img/p1_phy_s0_q06_3.png}{82.8}{0}
\qfigdef{p1_phy_s0_q06_4}{SRGBT - Test 06 (29 Aug 2021)-img/p1_phy_s0_q06_4.png}{83.1}{0}
\qfigdef{p1_phy_s0_q06_5}{SRGBT - Test 06 (29 Aug 2021)-img/p1_phy_s0_q06_5.png}{83.2}{0}
\qfigdef{p1_phy_s0_q07_1}{SRGBT - Test 06 (29 Aug 2021)-img/p1_phy_s0_q07_1.png}{106.5}{0}
\qfigdef{p1_phy_s0_q08_2}{SRGBT - Test 06 (29 Aug 2021)-img/p1_phy_s0_q08_2.png}{94.7}{0}
\qfigdef{p1_phy_s0_q09_1}{SRGBT - Test 06 (29 Aug 2021)-img/p1_phy_s0_q09_1.png}{94.3}{0}
\qfigdef{p1_phy_s0_q10_1}{SRGBT - Test 06 (29 Aug 2021)-img/p1_phy_s0_q10_1.png}{90.9}{0}
\qfigdef{p1_phy_s0_q10_2}{SRGBT - Test 06 (29 Aug 2021)-img/p1_phy_s0_q10_2.png}{270.6}{1}
\qfigdef{p1_phy_s0_q11_1}{SRGBT - Test 06 (29 Aug 2021)-img/p1_phy_s0_q11_1.png}{128.7}{0}
\qfigdef{p1_phy_s0_q12_1}{SRGBT - Test 06 (29 Aug 2021)-img/p1_phy_s0_q12_1.png}{177.9}{0}
\qfigdef{p1_phy_s0_q13_1}{SRGBT - Test 06 (29 Aug 2021)-img/p1_phy_s0_q13_1.png}{66.4}{0}
\qfigdef{p1_phy_s0_q14_1}{SRGBT - Test 06 (29 Aug 2021)-img/p1_phy_s0_q14_1.png}{68.2}{0}
\qfigdef{p1_phy_s0_q14_2}{SRGBT - Test 06 (29 Aug 2021)-img/p1_phy_s0_q14_2.png}{94.7}{0}
\qfigdef{p1_phy_s0_q16_1}{SRGBT - Test 06 (29 Aug 2021)-img/p1_phy_s0_q16_1.png}{108.7}{0}
\qfigdef{p1_phy_s0_q17_2}{SRGBT - Test 06 (29 Aug 2021)-img/p1_phy_s0_q17_2.png}{263.1}{1}
\qfigdef{p1_phy_s0_q18_1}{SRGBT - Test 06 (29 Aug 2021)-img/p1_phy_s0_q18_1.png}{97.4}{0}
\qfigdef{p1_phy_s0_q20_2}{SRGBT - Test 06 (29 Aug 2021)-img/p1_phy_s0_q20_2.png}{47.1}{0}
\qfigdef{p1_phy_s0_q21_1}{SRGBT - Test 06 (29 Aug 2021)-img/p1_phy_s0_q21_1.png}{313.2}{1}
\qfigdef{p1_phy_s0_q22_1}{SRGBT - Test 06 (29 Aug 2021)-img/p1_phy_s0_q22_1.png}{313.2}{1}
\qfigdef{p1_math_s0_pre_1}{SRGBT - Test 06 (29 Aug 2021)-img/p1_math_s0_pre_1.png}{83.4}{0}
\raggedcolumns
\makeatletter\@beginparpenalty=10000\makeatother
\newcolumntype{L}[1]{>{\raggedright\arraybackslash}p{#1}}
\begin{document}
\begin{center}{\Large\bfseries SRGBT --- Test 06 (29 Aug 2021)}\end{center}
\begin{multicols}{2}
\raggedcolumns
\papersec{1}
\partsec{1}{PHYSICS}
\qn{1} A gas is transferred from a state 1 to a state 2 in two ways. (1) Directly by an isobar, and (2) first by the isochor 1-3, then by the isobar 3-4, and finally, by the isochor 4-2. the entropy increment in both cases is $\Delta\mathrm S_{\text{process-1}}=\mathrm K\,\Delta\mathrm S_{\text{process-2}}$ the same. Find K.
\begin{qoptions}
\opt{A}{$\sqrt2$}
\opt{B}{$\sqrt3$}
\opt{C}{1}
\opt{D}{None of these}
\end{qoptions}
\qn{2} Two objects with initial temperatures T$_1$ and T$_2$ (with T$_1>$ T$_2$) are brought into contact. The objects are isolated from their surroundings, and the masses (m) and specific heat capacities (c) of the two objects coincide. How does the total entropy of these objects change as the temperatures become equal ?
\begin{qoptions}
\opt{A}{cm $\ln\left[\dfrac{(\mathrm T_1+\mathrm T_2)^2}{4\mathrm T_1\mathrm T_2}+1\right]$}
\opt{B}{cm $\ln\left[\dfrac{(\mathrm T_1-\mathrm T_2)^2}{4\mathrm T_1\mathrm T_2}+1\right]$}
\opt{C}{cm $\ln\left[\dfrac{(\mathrm T_1+\mathrm T_2)^2}{4\mathrm T_1\mathrm T_2}-1\right]$}
\opt{D}{None of these}
\end{qoptions}
\qn{3} 0.01 moles of an ideal diatomic gas is enclosed in an adiabatic cylinder of cross-sectional area A = 10$^{-4}$m$^2$. In the arrangement shown, a block of mass M = 0.8 kg is placed on a horizontal support, and piston of mass m = 1 kg is suspended from a spring of stiffness constant k = 16 N/m. Initially, the spring is relaxed and the volume of the gas is V = 1.4 $\times$ 10$^{-4}$ m$^3$.\par
\qfig{p1_phy_s0_q03_1}
When the gas in the cylinder is heated up the piston starts moving up and the spring gets compressed so that the block M is just lifted up. Determine the heat supplied (in Joule).\par
Take atmospheric pressure P$_0$ = 10$^5$ Nm$^{-2}$, g = 10 m/s$^2$.
\begin{qoptions}
\opt{A}{25}
\opt{B}{50}
\opt{C}{75}
\opt{D}{55}
\end{qoptions}
\qn{4} Two adiabatic processes involving an ideal gas are plotted on a P-V diagram. A and B are two points on these curves as shown in the diagram. Which of the following statements is true ?\par
\qfig{p1_phy_s0_q04_1}
\begin{qoptions}
\opt{A}{Heat is given to system in the process AB}
\opt{B}{Heat is rejected by system in the process AB}
\opt{C}{Whether heat is rejected or absorbed by the system during the process AB depends on the magnitude of the slope of line AB}
\opt{D}{Whether heat is rejected or absorbed by the system during the process AB depends on the position of the points A and B on the P-V diagram.}
\end{qoptions}
\qn{5} A movable heavy piston is supported by a spring inside a vertical cylindrical container as shown. When all air is pumped out of the container, the piston is in equilibrium as shown in Diagram 1, with only a tiny gap between the piston and the bottom of the container. When a portion of gas at temperature T is introduced under the piston, the later rises to a height h as shown in Diagram 2. What would be the height of the piston above the bottom of the container if the gas is then heated to a temperature 2T ? Assume that the piston moves without friction and that the spring obeys Hooke's law.\par
\qfig{p1_phy_s0_q05_1}
\begin{qoptions}
\opt{A}{$\sqrt2\,$h}
\opt{B}{$\sqrt3\,$h}
\opt{C}{$\sqrt4\,$h}
\opt{D}{$\sqrt5\,$h}
\end{qoptions}
\qn{6} One mole of a monoatomic gas undergoes the processes 1-2 and 2-3 as shown the corresponding graph for the above process for 2 moles of the gas is\par
\qfig{p1_phy_s0_q06_1}
\begin{qoptions}
\opt{A}{\qfigi{p1_phy_s0_q06_2}}
\opt{B}{\qfigi{p1_phy_s0_q06_3}}
\opt{C}{\qfigi{p1_phy_s0_q06_4}}
\opt{D}{\qfigi{p1_phy_s0_q06_5}}
\end{qoptions}
\qn{7} One mole of monoatomic gas is taken through cyclic process shown below. T$_\mathrm A$ = 300 K. Process AB is defined as PT = constant.\par
\qfig{p1_phy_s0_q07_1}
\begin{qoptions}
\opt{A}{Work done in process AB is $-400$ R.}
\opt{B}{Change in internal energy in process CA is 900 R.}
\opt{C}{Heat transferred in the process BC is 2000 R.}
\opt{D}{Change in internal energy in process CA is $-900$ R.}
\end{qoptions}
\qn{8} A sample of ideal gas of some mass is taken through the cyclic process shown in the figure. The temperature of the gas at state A is T$_\mathrm A$ = 200 K, and at state C is T$_\mathrm C$ = 1800 K. Choose the correct option.\par
\qfig{p1_phy_s0_q08_2}
\begin{qoptions}
\opt{A}{$\dfrac{\mathrm V_\mathrm A}{\mathrm V_\mathrm C}=\dfrac13$}
\opt{B}{$\dfrac{\mathrm V_\mathrm A}{\mathrm V_\mathrm B}=\dfrac19$}
\opt{C}{$\dfrac{\mathrm P_\mathrm A\mathrm V_\mathrm A}{\mathrm P_\mathrm C\mathrm V_\mathrm C}=\dfrac13$}
\opt{D}{$\dfrac{\mathrm P_\mathrm A\mathrm V_\mathrm A}{\mathrm P_\mathrm C\mathrm V_\mathrm C}=\dfrac19$}
\end{qoptions}
\qn{9} A sample of ideal gas is taken through the cyclic process shown in the figure. The temperature of the gas at A is T$_\mathrm A$ = 200 K. At states B and C the temperature of the gas is the same. Choose the correct option :\par
\qfig{p1_phy_s0_q09_1}
\begin{qoptions}
\opt{A}{P$_\mathrm B$ = 3P$_\mathrm C$}
\opt{B}{P$_\mathrm C$ = 3P$_\mathrm A$}
\opt{C}{P$_\mathrm B$ = 9P$_\mathrm A$}
\opt{D}{None of these}
\end{qoptions}
\qfullbegin
\qn{10} Density ($\rho$) versus temperature (T) graph of a thermodynamic cycle of an ideal gas is as shown. If BC and AD are the part of rectangular hyperbola then which of the following graphs will represent the same thermodynamic cycle?\par
\qfig{p1_phy_s0_q10_1}
\qfullsplit
\qfig{p1_phy_s0_q10_2}
\qfullend
\qn{11} Two cylinders having frictionless pistons A and B are connected by small tube through which a light inextensible string (which is connected to centres of both the pistons) passes. Both the pistons are having same thickness and made up of same material. Radius of piston A is twice the radius of piston B. P is the pressure of gas enclosed between the pistons and P$_0$ be the atmospheric pressure. String is in horizontal plane\par
\qfig[0.85]{p1_phy_s0_q11_1}
\begin{qoptions}
\opt{A}{If tension in the string is non zero both the pistons must accelerate}
\opt{B}{Both the pistons will be in equilibrium only when P = P$_0$}
\opt{C}{If string is in vertical plane such that the piston A is on the top, the equilibrium of pistons is not possible}
\opt{D}{If string is in vertical plane such that the piston B is on the top, the equilibrium of pistons is not possible}
\end{qoptions}
\qn{12} The temperature of a well stirred liquid kept open to a cold surrounding is plotted against time. The value of $\sec^2\theta_1$ is :\par
\qfig[0.8]{p1_phy_s0_q12_1}
\begin{qoptions}
\opt{A}{$1+9\tan^2\theta_2$}
\opt{B}{$1+\tan^2\theta_2$}
\opt{C}{$1+3\tan^2\theta_2$}
\opt{D}{$3+\tan^2\theta_2$}
\end{qoptions}
\qn{13} An ideal gas undergoes a cyclic process ABCA as shown in P-V diagram. The gas expands isothermally along AB and adiabatically along BC and the respective work done are 600 J and 250 J. If the efficiency of the cycle is 1/3 then\par
\qfig{p1_phy_s0_q13_1}
\begin{qoptions}
\opt{A}{during the compression, the gas gives up 400J of heat to the surrounding}
\opt{B}{the amount of work done on the gas along CA is 300 J}
\opt{C}{the net work done is 200 J}
\opt{D}{the change in the internal energy along CA is 250 J}
\end{qoptions}
\qn{14} An ideal gas undergoes an expansion from a state with initial temperature T$_1$ and initial volume V$_1$ to final volume V$_2$ through three different polytropic processes A, B and C as shown in the P-V diagram. If $|\Delta\mathrm E_\mathrm A|$, $|\Delta\mathrm E_\mathrm B|$ and $|\Delta\mathrm E_\mathrm C|$ be the changes in internal energy along the three paths respectively, then :-\par
\qfig{p1_phy_s0_q14_1}
\begin{qoptions}
\opt{A}{$|\Delta\mathrm E_\mathrm A|<|\Delta\mathrm E_\mathrm B|<|\Delta\mathrm E_\mathrm C|$ if temperature in every process decreases}
\opt{B}{$|\Delta\mathrm E_\mathrm A|>|\Delta\mathrm E_\mathrm B|>|\Delta\mathrm E_\mathrm C|$ if temperature in every process decreases}
\opt{C}{$|\Delta\mathrm E_\mathrm A|>|\Delta\mathrm E_\mathrm B|>|\Delta\mathrm E_\mathrm C|$ if temperature in every process increases}
\opt{D}{$|\Delta\mathrm E_\mathrm B|<|\Delta\mathrm E_\mathrm A|<|\Delta\mathrm E_\mathrm C|$ if temperature in every process increases}
\end{qoptions}
\paraq{15}{16}In the following P-V diagram, two adiabatics cut two isothermal at T$_1$ and T$_2$. 2 moles of an ideal diatomic gas is taken through the cyclic process ABCDA.\par
\qfig{p1_phy_s0_q14_2}
\qn{15} In which process heat is absorbed by the gas :-
\begin{qoptions}
\opt{A}{DA}
\opt{B}{BC}
\opt{C}{AB}
\opt{D}{CD}
\end{qoptions}
\qn{16} The volumes at different points A,B,C \& D are satisfy the following relation :-
\begin{qoptions}
\opt{A}{V$_\mathrm A$V$_\mathrm B$ = V$_\mathrm C$V$_\mathrm D$}
\opt{B}{V$_\mathrm A$V$_\mathrm C$= V$_\mathrm B$V$_\mathrm D$}
\opt{C}{V$_\mathrm A$V$_\mathrm D$ = V$_\mathrm B$V$_\mathrm C$}
\opt{D}{V$_\mathrm A$V$_\mathrm C^2$ = V$_\mathrm B$V$_\mathrm D^2$}
\end{qoptions}
\paraq{17}{18}A sample of mono-atomic ideal gas is taken through a cyclic process ABCDA. The cyclic process is such that each of the subprocesses have constant molar specific heat capacity. A graph is plotted in which heat exchanged by the gas upto an instant is shown on x-axis and work done by the gas upto that instant is shown on y-axis.\par
\qfig{p1_phy_s0_q16_1}
\qfullbegin
\qn{17} The given cyclic process can be represented on a pressure-volume graph as :-\par
\qfig{p1_phy_s0_q17_2}
\qfullend
\qn{18} Choose INCORRECT statement :-
\begin{qoptions}
\opt{A}{Magnitudes of temperature change in the process AB and CD are same}
\opt{B}{Process AB and CD are isochoric}
\opt{C}{Process BC and DA are isothermal}
\opt{D}{Process BC is adiabatic but CD is isothermal}
\end{qoptions}
\paraq{19}{20}A thermally insulated cylindrical container is divided into three parts by two insulated pistons which can move without friction. The central part of the container contains n moles of oxygen gas. The left and right parts of the container have vacuum. The natural length of both the spring is 2$\ell_0$. Both the springs are initially compressed by $\ell_0$. Heating of the gas is done with the help of a heating coil H such that pistons move slowly by distance x as shown in the figure. Neglect the heat capacity of the container, pistons and springs. Spring constant of both the spring is k.\par
\qfig{p1_phy_s0_q18_1}
Now answer the following questions
\qn{19} The change in internal energy of the gas during the given displacement x is
\begin{qoptions}
\opt{A}{5kx$^2$}
\opt{B}{10k$\ell_0$x}
\opt{C}{5k(2$\ell_0$x + x$^2$)}
\opt{D}{can't be calculated.}
\end{qoptions}
\qn{20} The molar heat capacity of the gas during the process is
\begin{qoptions}
\opt{A}{$\mathrm R\left(\dfrac{13}2\ell_0\mathrm x+4\mathrm x^2\right)\Big/\,3\ell_0\mathrm x+2\mathrm x^2$}
\opt{B}{2R}
\opt{C}{3R}
\opt{D}{$\mathrm R\left(\dfrac{15}2\ell_0\mathrm x+4\mathrm x^2\right)\Big/\,5\ell_0\mathrm x+2\mathrm x^2$}
\end{qoptions}
\paraq{21}{22}The figure represents a cyclic process in which the process CA is adiabatic\par
\qfig{p1_phy_s0_q20_2}
\qfullbegin
\qn{21} Choose the diagram showing the correct representation of the thermodynamic cycle on at P--T diagram\par
\qfig{p1_phy_s0_q21_1}
\qn{22} Choose the diagram showing the correct representation of the thermodynamic cycle on at V--T diagram\par
\qfig{p1_phy_s0_q22_1}
\qfullsplit
\raggedcolumns
\qfullend
\papersec{1}
\partsec{2}{CHEMISTRY}
\qn{23} Boric acids is a weak acid in aqueous solution. But its acidity increases significantly in the presence of ethylene glycol, because
\begin{qoptions}
\opt{A}{ethylene glycol releases additional H$^+$.}
\opt{B}{B(OH)$_4^-$ is consumed in forming a compound with ethylene glycol.}
\opt{C}{ethylene glycol neutralizes H$^+$ released by boric acid.}
\opt{D}{Boric acid dissociates better in the mixed solvent.}
\end{qoptions}
\qn{24} A soft metal (A) gives violet flame test. When burnt in excess O$_2$, (A) gives yellow powder (B) which on reaction with water gives alkaline solution C and O$_2$. Identify A and B.
\begin{qoptions}
\opt{A}{Na, Na$_2$O$_2$}
\opt{B}{K, K$_2$O$_2$}
\opt{C}{K, KO$_2$}
\opt{D}{Li, Li$_2$O$_2$}
\end{qoptions}
\qn{25} How many statements are correct.\par
(i) Ethyl borate, B(OC$_2$H$_5$)$_3$ burns with green edged flame.\par
(ii) Boron is obtained by the pyrolysis of BI$_3$.\par
(iii) One mole of borax reacts with one mole of acid.\par
(iv) In sodium peroxoborate, each boron is sp$^3$ hybridised.\par
(v) H$_3$BO$_3$ does not dissolve in aqueous HF.\par
(vi) B$_2$H$_6$ + NH$_3$ $\xrightarrow[\text{higher tempt.($>200^\circ$C)}]{\text{Excess NH}_3}$ B$_3$N$_3$H$_6$ (borazine).\par
(vii) The graphite is paramagnetic and diamond is diamagnetic in nature.\par
(viii) Graphite acts as a metallic conductor along the layers of carbon atoms and as semi-conductor perpendicular to the layers of carbon atoms.\par
(ix) Graphite is less denser than diamond.\par
(x) Silicones are strongly water repellent.
\begin{qoptions}
\opt{A}{3}
\opt{B}{5}
\opt{C}{6}
\opt{D}{8}
\end{qoptions}
\qn{26} Which of the following statements is \textbf{INCORRECT} w.r.t Boron for its compounds.
\begin{qoptions}
\opt{A}{Boric acid with NaOH does not give a sharp end point but when taken along with sugars it gives a sharp end point with NaOH.}
\opt{B}{Aqueous solution of Borax shows buffer properties.}
\opt{C}{B$_2$H$_6$ on reaction with excess NH$_3$ at low temperatures forms an ionic compound which on heating gives an aromatic compound which is less reactive than benzene.}
\opt{D}{On hydrolysis of BF$_3$, boric acid along with flouroborates are obtained.}
\end{qoptions}
\qn{27} Identify the correct statement :
\begin{qoptions}
\opt{A}{Gypsum contains a lower percentage of calcium than Plaster of Paris}
\opt{B}{Gypsum is obtained by heating Plaster of Paris}
\opt{C}{Plaster of Paris is obtained by hydration of gypsum}
\opt{D}{Plaster of Paris is obtained by partial oxidation of gypsum}
\end{qoptions}
\qn{28} Which is not correctly matched ?\par
(1) Basic strength of oxides : Cs$_2$O $<$ Rb$_2$O $>$ K$_2$O $<$ Na$_2$O\par
(2) Stability of peroxides : Na$_2$O$_2$ $<$ K$_2$O$_2$ $<$ Rb$_2$O$_2$ $<$ Cs$_2$O$_2$\par
(3) Stability of bicarbonates : LiHCO$_3$ $<$ NaHCO$_3$ $<$ KHCO$_3$ $<$ RbHCO$_3$\par
(4) Melting point : NaF $<$ NaCl $<$ NaBr $<$ NaI
\begin{qoptions}
\opt{A}{1 and 4}
\opt{B}{1 and 3}
\opt{C}{1 and 2}
\opt{D}{2 and 3}
\end{qoptions}
\qn{29} KO$_2$ produces observable change with
\begin{qoptions}
\opt{A}{Mg$^{2+}$ (aq)}
\opt{B}{Fe$^{2+}$ (aq)}
\opt{C}{Cr$^{3+}$ (aq)}
\opt{D}{Al$^{3+}$ (aq)}
\end{qoptions}
\qn{30} Which of the following statements are correct about the dilute solutions of alkali metals in liquid ammonia ?
\begin{qoptions}
\opt{A}{They are deep blue coloured solution}
\opt{B}{They are highly conducting in nature}
\opt{C}{They are diamagnetic in nature}
\opt{D}{Ammoniated cation and solvated electrons are formed in the solution.}
\end{qoptions}
\qn{31} Which of the following are suitable methods for the preparation of Li$_2$CO$_3$ ?
\begin{qoptions}
\opt{A}{Li$_2$O + CO$_2$ $\longrightarrow$}
\opt{B}{LiOH + CO$_2$ $\longrightarrow$}
\opt{C}{LiOH + H$_2$CO$_3$ $\longrightarrow$}
\opt{D}{Li$_2$O$_2$ + CO$_2$ $\longrightarrow$}
\end{qoptions}
\qn{32} Which of the following lewis bases form adduct through symmetrical cleavage of B$_2$H$_6$ ?
\begin{qoptions}
\opt{A}{Pyridine}
\opt{B}{CO}
\opt{C}{T.H.F.}
\opt{D}{O(Et)$_2$}
\end{qoptions}
\qn{33} With respect to graphite and diamond, which of the statement(s) given below is (are) \textbf{CORRECT} ?
\begin{qoptions}
\opt{A}{Graphite is harder than diamond.}
\opt{B}{Graphite has higher electrical conductivity than diamond.}
\opt{C}{Graphite has higher thermal conductivity than diamond.}
\opt{D}{Graphite has higher C -- C bond order than diamond.}
\end{qoptions}
\qn{34} Which of the following can not be converted into anhydrous chloride only by heating ?
\begin{qoptions}
\opt{A}{MgCl$_2$.6H$_2$O}
\opt{B}{FeCl$_3$.6H$_2$O}
\opt{C}{AlCl$_3$.6H$_2$O}
\opt{D}{CrCl$_3$.6H$_2$O}
\end{qoptions}
\qn{35} Carbon monoxide is not prepared by:
\begin{qoptions}
\opt{A}{Heating formic acid with conc. H$_2$SO$_4$.}
\opt{B}{Heating potassium ferrocyanide with conc.H$_2$SO$_4$.}
\opt{C}{Heating malonic acid with P$_4$O$_{10}$.}
\opt{D}{Hydrolysis of Mg$_2$C$_3$.}
\end{qoptions}
\qn{36} Which of the following is/are correct.
\begin{qoptions}
\opt{A}{Ag$_2$S is much less soluble than Ag$_2$O.}
\opt{B}{Fe(OH)$_3$ is much less soluble than Fe(OH)$_2$.}
\opt{C}{BeCO$_3$ is less soluble than MgCO$_3$.}
\opt{D}{Melting point of AlCl$_3$ is much less than that of NaCl.}
\end{qoptions}
\paraq{37}{38}\par
(i) Boron + O$_2$ $\xrightarrow{700^\circ\mathrm C}$ (X)\par
(ii) [X] + C(carbon) + Cl$_2$ $\longrightarrow$ (Y) + CO\par
(iii) [Y] + LiAlH$_4$ $\longrightarrow$ (Z) + LiCl + AlCl$_3$\par
(iv) [Z] + NH$_3$ $\longrightarrow$ (A) $\xrightarrow{\text{heat}}$ (B)\par
(v) [Z] + NaH $\longrightarrow$ (D)
\qn{37} Incorrect about compound (Z) is:
\begin{qoptions}
\opt{A}{A covalent compound}
\opt{B}{A electron deficient compound}
\opt{C}{3c--2e$^-$ compound}
\opt{D}{Having ethane like structure}
\end{qoptions}
\qn{38} Compound (B) is:
\begin{qoptions}
\opt{A}{Boron nitride}
\opt{B}{Borazole}
\opt{C}{Borazon}
\opt{D}{Organic benzene}
\end{qoptions}
\paraq{39}{40}Lithium shows diagonal relationship with Magnesium as well as lithium shows anomalous properties with other alkali metals.
\qn{39} Which of the following properties does not represent the diagonal relationship between them:-
\begin{qoptions}
\opt{A}{Li$_2$O and MgO both do not combine with excess oxygen to form superoxide}
\opt{B}{LiCl and MgCl$_2$ are soluble in ethanol}
\opt{C}{Li$^+$ forms complex with co-ordination number of 4 while Mg$^{2+}$ forms complex with co-ordination number of 6}
\opt{D}{LiCl and MgCl$_2$ both are deliquescent}
\end{qoptions}
\qn{40} Which of the following properties is wrong towards the anomalous properties of Li :-
\begin{qoptions}
\opt{A}{Lithium forms oxide as well as nitride when burns in air}
\opt{B}{LiHCO$_3$ does not exist in solid form}
\opt{C}{Melting point of Lithium is highest among alkali metals}
\opt{D}{Li is kinetically most reactive towards water}
\end{qoptions}
\paraq{41}{42}\par
[A] $\xrightarrow[\Delta]{\mathrm{Na_2CO_3(aq.)}}$ CaCO$_3$ (white ppt.) + B(aq.) + C(aq.)\par
[B]$_{\mathrm{(aq.)}}$ $\xrightarrow{\text{On crystallisation}}$ [X] Crystallised\par
[B]$_{\mathrm{(s)}}$ + CoO $\xrightarrow{\Delta}$ [Blue coloured bead (D)]
\qn{41} ``A'' is :
\begin{qoptions}
\opt{A}{CaCO$_3$}
\opt{B}{Ca$_2$B$_6$O$_{11}$.5H$_2$O}
\opt{C}{Co(BO$_2$)$_2$}
\opt{D}{Ca$_3$Si$_2$O$_7$}
\end{qoptions}
\qn{42} ``X'' is :
\begin{qoptions}
\opt{A}{NaBO$_2$}
\opt{B}{Na$_2$B$_4$O$_7$.10H$_2$O}
\opt{C}{Na$_2$B$_4$O$_7$.8H$_2$O}
\opt{D}{Na$_2$[B$_4$O$_5$(OH)$_4$]}
\end{qoptions}
\paraq{43}{44}When Pb(NO$_3$)$_2$ is heated then it forms two gaseous products \textbf{A} and \textbf{B}. When nitrogen containing gaseous product `\textbf{B}' is reacted with another neutral oxide of Nitrogen then it forms a coloured liquid \textbf{C} at low temperature. Non nitrogen containg gas `\textbf{A}' when heated in silent electric dischange then forms \textbf{D} which is strong bleaching agent. Then
\qn{43} Which is \textbf{not true} about compound `C'
\begin{qoptions}
\opt{A}{it contains 2 N-atoms}
\opt{B}{Both the N - atoms are in sp$^2$ hybridisation}
\opt{C}{It has 8 non bonded `lp' of e$^-$}
\opt{D}{All `N -- O' bond length are same}
\end{qoptions}
\qn{44} Identify true statement
\begin{qoptions}
\opt{A}{Both A and D are paramagnetic}
\opt{B}{In D all bond lengths are equal}
\opt{C}{entropy of system increases in the conversion of A to D}
\opt{D}{D is colourless in solid and liquid state.}
\end{qoptions}
\raggedcolumns
\papersec{1}
\partsec{3}{MATHEMATICS}
\qn{45}
\qfig{p1_math_s0_pre_1}
In how many ways 5X can be placed in the square of the figure given above so that no row remain empty is:
\begin{qoptions}
\opt{A}{126}
\opt{B}{118}
\opt{C}{98}
\opt{D}{80}
\end{qoptions}
\qn{46} If a person can scores 0,2,3 or 5 points in a match then total numbers of ways by which he can score 15 points in 10 matches is :
\begin{qoptions}
\opt{A}{7752}
\opt{B}{16752}
\opt{C}{14400}
\opt{D}{16500}
\end{qoptions}
\qn{47} Let A = \{x : 1 $\le$ x $\le$ 2021\}, y is any number belonging to set A such that total number of divisor of y is odd then number of such possible number y in set A is :
\begin{qoptions}
\opt{A}{1985}
\opt{B}{44}
\opt{C}{1986}
\opt{D}{45}
\end{qoptions}
\qn{48} Number of ways in which 7 different balls can be distributed among 4 children if each child receive atleast 1 ball is:
\begin{qoptions}
\opt{A}{5880}
\opt{B}{3360}
\opt{C}{8400}
\opt{D}{7560}
\end{qoptions}
\qn{49} Let xy = n, x, y $\in$ N where n is a factor of 120 then number of possible orders pair (x, y) is :
\begin{qoptions}
\opt{A}{16}
\opt{B}{90}
\opt{C}{64}
\opt{D}{32}
\end{qoptions}
\qn{50} Let x$_1$, x$_2$, x$_3$, \dots\dots, x$_k$ are the total divisors of positive integer N. if $\displaystyle\sum_{i=1}^{k}\mathrm x_i=360$ and $\displaystyle\sum_{i=1}^{k}\dfrac1{\mathrm x_i}=3$, then the value of k is
\begin{qoptions}
\opt{A}{4}
\opt{B}{8}
\opt{C}{15}
\opt{D}{16}
\end{qoptions}
\qn{51} Let sequence a$_1$, a$_2$, \dots\dots\dots, a$_{20}$ be some permutation of number 1,2,3, \dots\dots.20 such that, then number of such permutations, is
\begin{qoptions}
\opt{A}{Total number of permutation if 3 divides $|\mathrm a_i-i|\ \forall i=1,2,3,\dots\dots\dots18$, is $(7!)^3$}
\opt{B}{Total number of permutation if 4 divides $|\mathrm a_i-i|\ \forall i=1,2,3,\dots\dots\dots18$, is $(5!)^4$}
\opt{C}{Total number of permutation if 5 divides $|\mathrm a_i-i|\ \forall i=1,2,3\,,\dots\dots..18$ is $(4!)^5$}
\opt{D}{Total number of permutation if 3 divides $|\mathrm a_i-i|\ \forall i=1,2,3,\dots\dots.18$ is $(6!)(7!)^2$}
\end{qoptions}
\qn{52} Which of the following are correct?
\begin{qoptions}
\opt{A}{Total possible natural numbers which are smaller than $2\times10^7$ and are divisible by 3 and formed using the digits 0,1,2 is 1457}
\opt{B}{Total possible natural numbers which are smaller than $2\times10^6$ and are divisible by 3 and formed using the digits 0,1,2 is 486}
\opt{C}{Total possible positive palindrome which are less than 100 is 18}
\opt{D}{Total possible positive palindrome which are less than 1000 is 108}
\end{qoptions}
\qn{53} Which of the following statements are correct:
\begin{qoptions}
\opt{A}{Total number of 4 letter words which can be formed with the letter of the word INEFFECTIVE is 1422}
\opt{B}{Total number of ways of arrangements of 5 identical red ball and 5 identical green ball if no two green ball are adjacent is 6.}
\opt{C}{Total number of quadrilateral which can be formed with the vertices of 12 sided polygon if no side of quadrilateral and polygon coincides is 105}
\opt{D}{Number of ways in which 10 different balls can be distributed to 3 children if each receive atleast one is $3^{10}-3(2)^{10}+3$}
\end{qoptions}
\qn{54} Let a$_n$ represent the number of rectangles formed using the vertices of n sided regular polygon then
\begin{qoptions}
\opt{A}{a$_{12}$ = 15}
\opt{B}{a$_{15}$ = 0}
\opt{C}{a$_{18}$ = 36}
\opt{D}{a$_{12}$ = 12}
\end{qoptions}
\qn{55} Let A = \{ 1,2,3, \dots\dots.20\} x and y are two distinct numbers where x, y $\in$ A. Total number of possible ordered pair (x, y). If :
\begin{qoptions}
\opt{A}{(x +y) is divisible by 3 is 128}
\opt{B}{(x + y) is odd is 100}
\opt{C}{x.y is even is 145}
\opt{D}{(x+y) is divisible by 4 is 90}
\end{qoptions}
\qn{56} Let L$_1$, L$_2$, L$_3$ , L$_4$, L$_5$ and L$_6$ are 6 letters and E$_1$, E$_2$, E$_3$, E$_4$, E$_5$ and E$_6$ are corresponding 6 envelopes. Number of ways in which 6 letters can be placed in 6 envelopes if :
\begin{qoptions}
\opt{A}{No letter is placed in its correct envelope is 265}
\opt{B}{No letter is placed in correct envelope and L$_4$ is always placed in E$_3$ is 51}
\opt{C}{No letter is placed in correct envelope and L$_4$ is not placed in E$_2$ is 212}
\opt{D}{No letter is placed in corrected envelope and L$_1$ does not goes in E$_3$ and L$_2$ does not goes in E$_5$ is 173.}
\end{qoptions}
\qn{57} Let A = \{a$_1$, a$_2$, a$_3$\dots\dots.a$_{10}$\} be a set containing 10 elements. Two subsets P \& Q of it is formed independently.\par
L= Number of possible ways in which subsets can formed such that (P $-$ Q) contains exactly 2 elements\par
M = Number of possible ways in which subsets can formed such that (P $\cap$ Q) is Null set . Then
\begin{qoptions}
\opt{A}{L = $^{10}$C$_2$3$^8$}
\opt{B}{L = 3$^9$}
\opt{C}{M = 3$^{10}$}
\opt{D}{M = 2$^{10}$}
\end{qoptions}
\qn{58} Let f(x) = (x $-$ 1) (x $-$ 2) \dots.. (x $-$ 15) and N = f$'$(7) then
\begin{qoptions}
\opt{A}{Total number of divisor of N is 360}
\opt{B}{Let xy = N, x, y $\in$ N then number of possible ordered pair (x, y) is 180}
\opt{C}{Total numver of divisor of N which are divisible by 10 is 220}
\opt{D}{Total numver of odd ivisor of N is 30}
\end{qoptions}
\paraq{59}{60}In how many ways 30 Identical balls can be distributed to three child A, B and C.
\qn{59} If each child must receive atleast one ball is
\begin{qoptions}
\opt{A}{$^{32}$C$_2$}
\opt{B}{$^{31}$C$_2$}
\opt{C}{$^{30}$C$_2$}
\opt{D}{$^{29}$C$_2$}
\end{qoptions}
\qn{60} If number of balls received by A is greater than B is :
\begin{qoptions}
\opt{A}{496}
\opt{B}{240}
\opt{C}{248}
\opt{D}{406}
\end{qoptions}
\paraq{61}{62}Consider the word\par
MISSISSIPPI\par
How many different arrangement are there of all the letters of the word so that
\qn{61} No two I are consecutive :
\begin{qoptions}
\opt{A}{7350}
\opt{B}{1225}
\opt{C}{2450}
\opt{D}{3675}
\end{qoptions}
\qn{62} First `P' preceds `M' is :
\begin{qoptions}
\opt{A}{$^8$C$_4\times{}^{11}$C$_3$}
\opt{B}{$^8$C$_4\times{}^{11}$C$_3\times4$}
\opt{C}{$^8$P$_4\times{}^{11}$P$_3$}
\opt{D}{$^8$C$_4\times{}^{11}$C$_3\times2$}
\end{qoptions}
\paraq{63}{64}Consider 6 children C$_1$, C$_2$, C$_3$, C$_4$, C$_5$ and C$_6$ of different heights, child with higher suffix is taller.
\qn{63} Number of ways in which these six children can line up in a single row so that none of them is standing between the two children taller than him, is
\begin{qoptions}
\opt{A}{32}
\opt{B}{30}
\opt{C}{15}
\opt{D}{20}
\end{qoptions}
\qn{64} Number of ways in which 10 alike marbles can be distributed among them if every child gets atleast one marble but not more than 4, is
\begin{qoptions}
\opt{A}{60}
\opt{B}{30}
\opt{C}{90}
\opt{D}{120}
\end{qoptions}
\paraq{65}{66}Consider the word ``BREAKAGE''. The number of ways in which the letters of this word can be arranged such that
\qn{65} The two A's are not together is :
\begin{qoptions}
\opt{A}{7560}
\opt{B}{5760}
\opt{C}{5040}
\opt{D}{15120}
\end{qoptions}
\qn{66} Neither two A's nor two E's are together is :
\begin{qoptions}
\opt{A}{1800}
\opt{B}{5760}
\opt{C}{720}
\opt{D}{7560}
\end{qoptions}
\end{multicols}
\end{document}
